\section{\textbf{Knowledge Invention: The Defining Capability of Superintelligence}}
\label{sec:knowledge_invention}

Sections~\ref{sec:knowledge_fusion} and~\ref{sec:knowledge_discovery}
identified two pathways beyond existing knowledge. Fusion constructs a
consequential relation among previously separated bodies of knowledge.
Discovery establishes a previously unknown phenomenon, entity, mechanism,
variable, regularity, law, or formal object. Neither pathway alone fully
specifies the transition to superintelligence. A fused relation may remain an
untested hypothesis, and a discovered phenomenon may remain unexplained.
Knowledge invention occurs when a generally intelligent machine originates the
epistemic structure through which such a result becomes consequential,
validated, and incorporable into collective knowledge.

The term \emph{invention} is used deliberately. Natural phenomena are not
invented by the system, but theories, explanatory structures, representations,
methods, algorithms, architectures, and research programs are constructed.
Machine knowledge invention may therefore begin with discovery, fusion, or
their interaction, but it culminates in a new epistemic object that did not
previously exist within either collective human knowledge or the recognized
human research agenda.

\subsection{\textbf{Novelty, Usefulness, and Epistemic Validity}}

Research on creativity commonly treats originality and effectiveness as joint
requirements \cite{runco2012creativity}. Work on artificial creativity likewise
distinguishes novel combination, exploration within a conceptual space, and
transformation of that space \cite{boden1998creativity}. These distinctions are
relevant to knowledge invention but remain insufficient for scientific
attribution. An idea may be original and useful without being true, and a true
statement may be trivial, previously known, or accidentally generated.

Let $\kappa$ be a candidate epistemic contribution. Knowledge invention
requires at least four properties:

\begin{equation}
\begin{aligned}
N(\kappa)&=1,\\
C(\kappa)&=1,\\
V(\kappa)&=1,\\
O_x(\kappa)&=1,
\end{aligned}
\label{eq:invention_properties}
\end{equation}

where $N$ denotes novelty relative to collective human knowledge, $C$
consequentiality, $V$ domain-appropriate validation, and $O_x$ causal
origination by machine $x$. The conjunction is essential. Novelty without
consequence is arbitrary variation; consequence without validity is an
attractive error; validity without novelty is rediscovery; and novelty,
consequence, and validity without machine origination describe a human
invention assisted by a machine rather than machine knowledge invention.

The standard is therefore stronger than computational creativity. A system may
generate a novel analogy, design, proof sketch, molecule, or theory candidate
and still fail to invent knowledge. The output becomes knowledge only after its
claims survive evaluation independent of the originating process.

\subsection{\textbf{Discovery, Fusion, and Invention}}

Fusion and discovery are neither mutually exclusive nor guaranteed to produce
invention. Fusion may expose a contradiction that becomes a new research
question. Discovery may identify the missing entity needed to reconcile two
theories. A discovered mechanism may then enable further fusion across domains.
The interaction is recurrent:

\begin{equation}
\begin{aligned}
\operatorname{Fusion}_{x}
&\longrightarrow q_F
\longrightarrow \operatorname{Discovery}_{x},\\
\operatorname{Discovery}_{x}
&\longrightarrow \kappa_D
\longrightarrow \operatorname{Fusion}_{x},\\
\operatorname{Fusion}_{x}oplus
\operatorname{Discovery}_{x}
&\longrightarrow \kappa_I,
\end{aligned}
\label{eq:fusion_discovery_invention}
\end{equation}

where $q_F$ is a question exposed through fusion, $\kappa_D$ is a discovered
object, $\oplus$ denotes epistemic interaction rather than numerical addition,
and $\kappa_I$ is a candidate knowledge invention.

Existing systems demonstrate portions of this cycle. Literature-based
discovery can expose implicit connections among separated findings
\cite{swanson1986fish}; symbolic-regression systems can construct interpretable
relations from data \cite{schmidt2009laws,udrescu2020aifeynman}; and autonomous
laboratories can connect planning, experimentation, analysis, and revision
\cite{boiko2023autonomous,szymanski2023autonomous}. Algorithm-discovery systems
have also produced provably correct or operationally validated procedures
beyond prior human designs \cite{fawzi2022alphatensor,mankowitz2023alphadev}.
These results establish increasingly powerful components of machine invention.
The unknown-unknown criterion adds the requirement that the machine originate
the consequential object of inquiry rather than receive it as a predefined
optimization target.

\subsection{\textbf{Degrees of Machine Contribution}}

The presence of an AI system in a scientific workflow does not determine the
degree of its contribution. We distinguish four levels.

\emph{Machine-assisted invention} occurs when humans formulate the problem,
construct the decisive hypothesis or concept, and use AI for retrieval,
calculation, prediction, optimization, or execution. \emph{Machine-generated
candidate invention} occurs when a system proposes a novel hypothesis, theory,
method, or design, but novelty, validity, consequence, or attribution remains
unestablished. \emph{Machine-originated invention} occurs when the machine
makes the causally essential conceptual contribution and the candidate is
beyond the known human frontier, but validation is incomplete.
\emph{Validated machine knowledge invention} occurs only when all attribution
and epistemic requirements are satisfied.

These levels can be represented as

\begin{equation}
\begin{aligned}
L_1&=\operatorname{Assist}_{x},\\
L_2&=\operatorname{Generate}_{x}(\kappa),\\
L_3&=L_2\land N(\kappa)\land O_x(\kappa),\\
L_4&=L_3\land C(\kappa)\land V(\kappa).
\end{aligned}
\label{eq:contribution_levels}
\end{equation}

Only $L_4$ constitutes knowledge invention in the strong sense used by this
paper. The hierarchy prevents a fluent output from being promoted directly to
knowledge and prevents a successful human--AI collaboration from being
misrepresented as autonomous machine origination. It also permits credit to be
assigned without requiring an all-or-nothing judgment: a system may make an
important machine-originated contribution even while humans provide
experimental infrastructure, independent review, or downstream interpretation.

\subsection{\textbf{The Attribution Standard}}

Attribution requires both causal and epistemic evidence. Causally, the machine
must perform an operation without which the contribution would not have arisen
through the same workflow. Epistemically, that operation must supply the novel
relation, object, explanation, or method rather than only supporting work.

Let $W$ be the complete invention workflow, $H$ its human contributions, $x$
the candidate machine inventor, and $\kappa_I$ the resulting contribution. We
define the attribution score conceptually as

\begin{equation}
\begin{aligned}
\mathcal{A}_{x}(\kappa_I\mid W)
=\Psi(&\operatorname{CounterfactualNecessity}_{x},\\
&\operatorname{EpistemicCentrality}_{x},\\
&\operatorname{HumanSpecification}^{-1},\\
&\operatorname{Traceability}_{x}),
\end{aligned}
\label{eq:attribution_score}
\end{equation}

where counterfactual necessity measures whether removing the machine's
contribution prevents the invention, epistemic centrality measures whether the
machine supplied the decisive intellectual content, inverse human specification
reduces attribution when humans predetermined the result, and traceability
measures whether the machine's contribution can be reconstructed from records.
Equation~\eqref{eq:attribution_score} does not prescribe a universal numerical
metric. It identifies the evidence that an attribution procedure must assess.

The human-specification term is particularly important. If researchers define
the target, enumerate the candidate space, specify the evaluation function, and
interpret the winning result, a machine may discover a powerful solution while
remaining inside a human-created epistemic frame. This can still be a genuine
machine discovery, as demonstrated by algorithm-search systems
\cite{fawzi2022alphatensor,mankowitz2023alphadev}, but it does not by itself
satisfy the unknown-unknown requirement for superintelligence. Strong knowledge
invention requires the machine to participate causally in constructing or
transforming the frame itself.

\subsection{\textbf{The Validation Boundary}}

Before validation, even an extraordinary machine-originated output remains a
candidate. The distinction is

\begin{equation}
\begin{aligned}
\operatorname{Candidate}_{x}(\kappa_I)
&\land
\operatorname{IndependentValidate}(\kappa_I)=1\\
&\Longrightarrow
\operatorname{KnowledgeInvention}_{x}(\kappa_I).
\end{aligned}
\label{eq:validation_boundary}
\end{equation}

Independent validation does not require that every evaluator be human. It
requires that acceptance not depend solely on the originating model, the same
unexamined data, or a criterion constructed after observing the result. The
validator may be a formal proof checker, replicated experiment, independent
model, external laboratory, engineering deployment, or combination of methods.
The appropriate standard depends on the epistemic object.

Reproducibility and transparent reporting are central to reliable knowledge
accumulation \cite{munafo2017reproducible}. For machine knowledge invention,
the validation record must additionally preserve model versions, prompts or
policies, tools, data provenance, intermediate hypotheses, rejected
alternatives, experimental decisions, and human interventions. Without this
record, neither novelty nor causal attribution can be assessed reliably.

Validation may also fail. A candidate invention can be rejected, restricted to
a smaller domain, or revised into a different claim. Such failure is not
evidence that the system lacks intelligence. It is evidence that candidate
generation and knowledge invention must remain distinct, just as hypotheses
and established scientific knowledge remain distinct in human research.

\subsection{\textbf{Knowledge-Invention Proposition}}

The paper's central proposition can now be stated.

\emph{Knowledge-Invention Proposition:} An artificial system demonstrates
superintelligence when, as a generally intelligent agent, it causally
originates a consequential epistemic contribution beyond both collective human
knowledge and the recognized human research agenda, and that contribution
survives independent, domain-appropriate validation.

Let $\mathcal{K}_{H,t}$ denote collective human knowledge,
$\mathcal{Q}_{H,t}$ the recognized human research agenda, and $\kappa_I$ the
candidate contribution. Then

\begin{equation}
\begin{aligned}
\operatorname{SI}(x)
\Longleftrightarrow{}&
\operatorname{AGI}(x)\\
&\land \exists\kappa_I:
\kappa_I\notin\mathcal{K}_{H,t}\\
&\land q(\kappa_I)\notin\mathcal{Q}_{H,t}\\
&\land O_x(\kappa_I)=1\\
&\land C(\kappa_I)=1\\
&\land V(\kappa_I)=1\\
&\land
\mathcal{K}_{H,t+1}
=\mathcal{K}_{H,t}\cup\{\kappa_I\}.
\end{aligned}
\label{eq:knowledge_invention_proposition}
\end{equation}

The proposition preserves the comparative meaning of superintelligence as
capability beyond humans \cite{bostrom2014superintelligence}, but specifies
where the relevant boundary lies. The machine need not outperform every human
on every activity. It must perform an epistemic act beyond the collective human
frontier: recognizing knowledge humanity did not know it was missing and
originating a validated contribution that supplies it.

\subsection{\textbf{From Isolated Success to Demonstrated Capability}}

One accepted result may be insufficient to establish a stable capability for
knowledge invention. The contribution may result from chance, training-data
leakage, incomplete novelty search, unrecorded human guidance, or an evaluation
that favored the generated output. Superintelligence is a capability claim,
not merely a historical label assigned after one surprising event.

Let $\{\kappa_1,\ldots,\kappa_m\}$ be independently evaluated machine-originated
contributions across contexts $\{c_1,\ldots,c_m\}$. A stronger demonstration
requires

\begin{equation}
\begin{aligned}
\sum_{i=1}^{m}
\mathbf{1}
\bigl[\operatorname{KI}_{x}(\kappa_i)=1\bigr]
&\geq r,\\
\operatorname{Diversity}(c_1,\ldots,c_m)
&\geq \eta,\\
\operatorname{Replication}(\kappa_i)
&=1
\quad \text{for accepted } i,
\end{aligned}
\label{eq:repeatable_invention}
\end{equation}

where $r$ is a required number of successful contributions and $\eta$ a
context-diversity requirement. Their values should be set by an evaluation
protocol rather than chosen retrospectively. Repeated success across genuinely
different unknown unknowns would provide stronger evidence that the system
possesses a transferable invention capability rather than one specialized
search procedure.

The evidentiary record for each contribution should include source provenance,
novelty analysis, machine and human action logs, derivation history,
alternative hypotheses, predictions, experiments or proofs, independent
replication, failure conditions, and the final attribution decision. Such
records make the claim falsifiable: evidence of prior human specification,
unrecognized source overlap, failed replication, or absent causal contribution
can defeat the attribution.

Knowledge invention therefore defines the transition proposed in this paper.
Superintelligence is not determined by how much knowledge a machine contains,
how fluently it describes that knowledge, or how effectively it optimizes a
known task. It is determined by whether a generally intelligent machine can
originate knowledge humanity did not know it was missing. The next section
addresses the resulting methodological challenge: how to evaluate and falsify
this capability without specifying the unknown unknowns in advance.
